\documentclass[12pt]{article}
\usepackage[utf8]{inputenc}
\usepackage[T1]{fontenc}
\usepackage{amsmath}
\usepackage{amssymb}
\usepackage{amsthm}
\usepackage{geometry}
\geometry{margin=1in}
\usepackage{setspace}
\usepackage{hyperref}
\usepackage{booktabs}
\usepackage{graphicx}
\usepackage{float}
\usepackage{enumitem}
\doublespacing

\newtheorem{proposition}{Proposition}
\newtheorem{definition}{Definition}

% Math macros
\newcommand{\kB}{k_{\mathrm{B}}}
\newcommand{\Teff}{T_{\mathrm{eff}}}
\newcommand{\Deff}{D_{\mathrm{eff}}}

\title{Coherence time in biological oscillator assemblies bounds the rate of state registration}

\author{Ian Todd\\
Sydney Medical School, University of Sydney\\
Sydney, NSW, Australia\\
\texttt{itod2305@uni.sydney.edu.au}}

\date{\today}

\begin{document}

\maketitle

\begin{abstract}
Cognitive and biological processing depends on coherent oscillations across distributed assemblies. Yet the physical principles determining how fast these assemblies can synchronize---and thus how fast neural processes associated with cognition can proceed---remain incompletely formalized. We derive a quantitative framework showing that coherence time in coupled oscillator networks scales exponentially with coordination depth. For $M$ semi-independent modules requiring phase alignment within tolerance $\varepsilon$ at coherence $r$, the minimum interval between irreversible state registrations (``commits'') is dominated by rare-event first-passage time into low-measure alignment sets on the configuration manifold. This produces a fundamental speed-flexibility trade-off: increasing coordination depth $M$ expands combinatorial flexibility but slows commits exponentially; increasing coherence $r$ accelerates synchronization but restricts dynamics to low-dimensional attractors. We situate this coordination limit within a hierarchy of physical bounds---quantum speed limits, Landauer-Margolus-Levitin thermodynamic constraints, and energy-time uncertainty relations---showing that in biological neural systems, these quantum and thermodynamic bounds are typically non-binding; the dominant constraint is the time required for distributed oscillatory assemblies to achieve sufficient phase coherence. While neural oscillators provide the most empirically accessible measurements, the framework applies to any biological system whose function requires coordinated phase relations and sparse commitment events, from chromosomal segregation to circadian networks. We provide parameter estimation protocols and testable predictions including alpha entrainment effects on temporal acuity and dual-loop dissociations under arousal.
\end{abstract}

\textbf{Keywords:} biological oscillations, coherence time, coordination depth, temporal resolution, quantum speed limits, Landauer bound, coupled oscillators

\section{Introduction}

\subsection{Processing speed in biological oscillator assemblies}

A fundamental insight from systems neuroscience is that cognition emerges from coherent oscillations across neural assemblies, not merely from individual spike rates. Working memory maintenance involves theta-gamma phase-amplitude coupling organized in discrete oscillatory bursts~\cite{miller2018,lisman2013,lundqvist2016}. Attention selectively enhances inter-areal synchronization in gamma band~\cite{fries2015,womelsdorf2007}. Perceptual binding depends on transient phase alignment across visual cortices~\cite{singer1999,engel2001}. Long-range communication occurs preferentially during coherent states~\cite{varela2001}. The ``communication through coherence'' framework~\cite{fries2005} has become central to understanding neural computation.

Yet despite extensive empirical characterization, the physical principles governing \emph{how fast} distributed assemblies can achieve coherence---and thus how quickly the neural processes underlying cognitive operations can proceed---remain incompletely formalized. Why does perceptual binding require 30--50 ms rather than 3 ms or 300 ms? Why do larger assemblies integrating more information process more slowly? What determines the relationship between oscillation frequency and temporal acuity?

\subsection{Coherence time as a physical bound on processing rate}

Recent work argued that living intelligence is an emergent property of high-dimensional, self-sustaining dynamical systems that maintain coherence against decoherence~\cite{todd2025intelligence}, building on analyses of falsifiability limits in adaptive biological systems~\cite{todd2025biosystems} and timing inaccessibility in sub-Landauer computation~\cite{todd2025demon}. The present manuscript addresses the complementary question: given such a coherent high-dimensional system, what sets the fastest timescale on which it can register change and update its state?

We formalize this as the minimum interval between irreversible ``commit'' events---transitions in which distributed dynamics become readable to downstream mechanisms, consistent with the emerging view that perception operates in discrete cycles~\cite{vanrullen2016}. This commit interval is bounded by a coordination term (coherence time as a first-passage problem on a product manifold) together with quantum, noise, and power limits.

\textbf{Operational definition of commits in neural systems.} A commit is operationally detected when distributed oscillatory dynamics cross a threshold that enables downstream readout. Empirical signatures include: (i) phase-locking index exceeding threshold across specified modules for a minimum duration, (ii) onset of stable cross-frequency coupling (e.g., theta-gamma PAC), (iii) classifier-detectable state transition in downstream neural populations. We do not claim all neural computation reduces to discrete commits; rather, we claim that the \emph{fastest update of downstream state} is bounded by the time required for multi-module phase alignment. The framework addresses temporal resolution limits, not the full richness of neural dynamics.

\begin{definition}[Commit event]\label{def:commit}
Given $M$ modules with aggregate phases $\psi_1(t), \ldots, \psi_M(t)$ and within-module coherences $r_{\mathrm{intra},m}(t)$, a \textbf{commit} occurs at time $t^*$ if:
\begin{enumerate}
\item All modules are internally coherent: $r_{\mathrm{intra},m}(t^*) \geq r_{\mathrm{thr}}$ for all $m$.
\item Phases are mutually aligned: $\max_m |\psi_m(t^*) - \bar{\psi}(t^*)| \leq \varepsilon/2$, where $\bar{\psi}$ is the circular mean.
\item The alignment persists for dwell time $T_{\mathrm{dwell}}$.
\end{enumerate}
The commit is defined \emph{relative to a downstream readout mechanism}; different readouts may have different $(r_{\mathrm{thr}}, \varepsilon, T_{\mathrm{dwell}})$ parameters.
\end{definition}

\medskip
\noindent\fbox{\parbox{\dimexpr\linewidth-2\fboxsep-2\fboxrule}{%
\textbf{What commit events are (and aren't).}
This framework does \textbf{not} claim cognition is discrete or that neural activity has
sharp ``commits.'' It claims the \textbf{fastest reliable downstream readout} of multi-module
coordination is lower-bounded by the timescale of phase alignment.

\textbf{Empirical proxies for commit detection:}
\begin{itemize}[noitemsep]
\item \textbf{Phase-locking index (PLI)} threshold crossing: sustained $r_{\mathrm{inter}} > r_{\mathrm{thresh}}$ for dwell time $t_{\mathrm{dwell}}$
\item \textbf{Phase-amplitude coupling (PAC)} onset: stable cross-frequency coupling emerges
\item \textbf{Classifier-detectable state change}: e.g., decoding accuracy crosses threshold
\end{itemize}

These are \textbf{measurement definitions}, not ontological claims. Different tasks may
require different readout criteria.}}
\medskip

\noindent\fbox{\parbox{\dimexpr\linewidth-2\fboxsep-2\fboxrule}{%
\textbf{Minimal empirical pipeline.}
To estimate coherence time from neural data:
\begin{enumerate}[noitemsep]
\item Identify modules via connectivity/functional clustering $\to$ count $M$
\item Extract module-level phases $\psi_m(t)$ via Hilbert transform or band-pass filtering
\item Compute intra- and inter-module coherence: $r_{\mathrm{intra}}, r_{\mathrm{inter}}$
\item Define alignment threshold $\varepsilon$ from task requirements (e.g., decoding accuracy)
\item Estimate attempt rate $\lambda_{\mathrm{attempt}}$ from decorrelation time of $r_{\mathrm{inter}}(t)$
\item Measure first-passage times to sustained alignment $\to$ fit Eq.~\ref{eq:coh-time}
\end{enumerate}}}

\textbf{Definition (Ontological dimensionality).} Throughout this work, ``dimensionality'' refers to a physical property of the system, not a mathematical measure. Ontological dimensionality is a quantitative measure of effective connectivity: how many independent directions of coordinated variation are accessible to a system given its interaction structure. High-dimensional states correspond to weakly constrained dynamics in which many components vary semi-independently; low-dimensional states correspond to strongly coupled dynamics restricted to a small number of collective modes. Dimensionality in this sense is determined not by component count, but by the rank of the system's interaction structure. Different measurement schemes (participation ratios, manifold dimension, information-geometric metrics) provide projections of this underlying physical property, but do not constitute it~\cite{amari2016}.

This ontological dimensionality is state-dependent: the same physical system may occupy high- or low-dimensional regimes depending on coupling strength and task demands. Coherence time depends on ontological dimensionality---specifically, how many semi-independent constraints must be simultaneously satisfied for a commit to occur---not on any particular estimator.

\textbf{Operationalizing dimensionality.} While ontological dimensionality is a physical property, it must be estimated from measurable quantities. For the coherence time framework, we operationalize effective dimensionality as:
\begin{enumerate}
\item \textbf{Coordination depth $M$}: Count of semi-independent modules (functional clusters) whose order parameters must jointly exceed threshold. Estimated via community detection on band-limited coherence networks or anatomical parcellation.
\item \textbf{Effective independence $\alpha$}: Fraction of potential constraint dimensions that are actually independent. Estimated via:
\begin{itemize}
\item Participation ratio of module-phase covariance (Eq.~\ref{eq:deff-estimator}),
\item or empirical scaling of $\tau$ with $M$ across task conditions.
\end{itemize}
\item \textbf{Effective constraint dimension}: $d_{\mathrm{eff}} = \alpha(M-1)$ represents the number of semi-independent coordination requirements for a commit.
\end{enumerate}
These are \emph{estimators} of the underlying physical property. Future work may employ other proxies (information-geometric metrics, topological data analysis, dynamical systems dimension estimates), but the coherence time itself depends on the physical constraint structure, not the measurement method.

\subsection{Structure of this work}

We first derive the coherence time formula from principles of phase synchronization in coupled oscillator networks (\S2.1), then embed it in a unified temporal resolution bound combining quantum, noise, and power constraints (\S2.2). Section 2.3 situates our framework explicitly within the quantum and thermodynamic limits of computation, addressing the Landauer bound, Margolus-Levitin quantum speed limit, and energy-time uncertainty relation. Section 3 applies the framework to neural systems, demonstrating quantitative agreement with perceptual binding windows, tachypsychia dissociations, and metabolic scaling. Section 4 discusses the speed-flexibility trade-off, delay embedding as a preservation mechanism, generalization beyond neural systems, and testable predictions.

\section{Theory and Methods}

\subsection{The unified temporal resolution bound}

We model biological temporal processing as a sequence of \emph{commits}---thermodynamically irreversible events that register high-dimensional internal state as low-dimensional output. A commit requires: (1) dimensional collapse from distributed state ($\Deff$ dimensions) to registered outcome (few dimensions), (2) dissipation of at least $\kB T \ln 2$ per bit (Landauer bound), and (3) creation of a persistent measurement influencing future dynamics.

The minimum time between commits is bounded below by four physical constraints:

\begin{equation}\label{eq:main-bound}
\boxed{\tau_{\mathrm{eff}} \;\gtrsim\; \max\left[\tau_{\mathrm{QSL}},\; \tau_{\mathrm{SNR}},\; \tau_{\mathrm{coh}},\; \tau_{\mathrm{power}}\right]}
\end{equation}

This is a \textbf{lower bound}: the effective commit interval must be at least as long as the slowest constraint. The max operation is appropriate when these constraints act on largely independent subsystems that can be satisfied in parallel (e.g., quantum state evolution, signal detection, and phase coordination proceed concurrently). If constraints were strictly sequential (e.g., signal detection \emph{then} coordination \emph{then} dissipation), a sum would be more appropriate. In biological neural systems, the constraints are largely parallel, and the coherence time dominates by orders of magnitude, so the distinction has minimal practical impact.

\textbf{Epistemic status of each term.} We distinguish:
\begin{itemize}
\item \textbf{Hard physical bounds} ($\tau_{\mathrm{QSL}}$, Landauer energy-per-bit): These follow from fundamental physics and cannot be violated regardless of engineering or biological design.
\item \textbf{Model-based lower bounds} ($\tau_{\mathrm{SNR}}$): These assume specific noise models (AWGN, matched filtering); different noise structures or detection strategies could yield different bounds, though the order of magnitude is robust.
\item \textbf{Coordination-time estimates} ($\tau_{\mathrm{coh}}$): These are asymptotic rare-event estimates whose prefactor depends on mixing assumptions and correlation structure. The scaling exponent $(M-1)$ is robust; the prefactor requires empirical calibration.
\item \textbf{Power constraints} ($\tau_{\mathrm{power}}$): Model-dependent, assuming a specific relationship between dimensional collapse and energy cost.
\end{itemize}
The key claim is not that Eq.~\ref{eq:main-bound} is a theorem, but that $\tau_{\mathrm{coh}}$ dominates by orders of magnitude in biological systems, making the precise form of other constraints largely irrelevant for practical predictions.

The four terms are:

\paragraph{Quantum speed limits ($\tau_{\mathrm{QSL}}$).} The Mandelstam-Tamm and Margolus-Levitin bounds~\cite{mandelstam1945,margolus1998} establish
\[
\tau_{\mathrm{QSL}} = \max\left[\frac{\hbar}{2\Delta E},\; \frac{\pi\hbar}{2\langle E\rangle}\right]
\]
For biological temperatures ($\Delta E \sim \kB T$), $\tau_{\mathrm{QSL}} \sim 10^{-13}$ s.

\paragraph{Signal-to-noise limit ($\tau_{\mathrm{SNR}}$).} Matched-filter detection in additive white Gaussian noise (two-sided power spectral density $N_0$, units: power/Hz) over bandwidth $B$ achieves $\mathrm{SNR}(T) = 2P_s T/N_0$ where $P_s$ is signal power and $T$ is integration time~\cite{poor1994}. Reaching detection threshold $\rho_{\min} \sim 5$--10 requires
\[
\tau_{\mathrm{SNR}} = \frac{\rho_{\min} N_0}{2P_s}
\]

\paragraph{Power limit ($\tau_{\mathrm{power}}$).} Dimensional collapse by $\Delta\mathcal{D}$ at dimensional chemical potential $\lambda_{\mathcal{D}} \sim \kB\Teff$ costs energy $\sim \lambda_{\mathcal{D}}\Delta\mathcal{D}$. At sustained power $P$, commit rate is bounded by
\[
\tau_{\mathrm{power}} = \frac{\lambda_{\mathcal{D}}\Delta\mathcal{D}}{P}
\]

\paragraph{Coherence time ($\tau_{\mathrm{coh}}$).}
We model commits as threshold crossings requiring $M$ semi-independent modules to achieve simultaneous phase alignment within tolerance $\varepsilon$, building on the theory of coupled oscillators~\cite{kuramoto1984,acebron2005}.
Let $r$ denote the \textbf{inter-module coherence} $r_{\mathrm{inter}} = |M^{-1}\sum_m e^{i\psi_m}|$ where $\psi_m$ is the aggregate phase of module $m$. (We reserve $r_{\mathrm{global}}$ for the all-oscillator order parameter, which serves as a diagnostic but does not appear in Eq.~\ref{eq:coh-time}.) The von Mises concentration $\kappa(r)$ satisfies $r=I_1(\kappa)/I_0(\kappa)$~\cite{mardia2000}.

\textbf{Why von Mises?} Given circular symmetry and a fixed mean resultant length $r$, the maximum-entropy distribution on the circle is von Mises~\cite{mardia2000}. This makes our modeling choice the canonical least-informative closure: $r \to \kappa \to p_1$ is uniquely determined without additional distributional assumptions. We validate this closure empirically in simulation (conditional-on-$r$ fitting yields $\kappa_{\mathrm{fitted}}/\kappa_{\mathrm{expected}} = 1.06 \pm 0.12$).

\textbf{Definition of alignment event.} We define alignment relative to the \textbf{instantaneous mean phase} $\bar{\psi} = \arg(\sum_m e^{i\psi_m})$: all modules must fall within $\varepsilon/2$ of $\bar{\psi}$. Formally:
\[
A_\varepsilon = \{|\psi_m - \bar{\psi}| < \varepsilon/2 \text{ for all } m\}
\]
This definition is natural for empirical detection (threshold on max phase deviation from mean) and matches our simulation event detector. Because the mean phase $\bar{\psi}$ is determined by the configuration itself, we have $M-1$ independent constraints (the mean absorbs one degree of freedom).

The alignment probability is:
\[
P(A_\varepsilon) \approx p_1(\varepsilon, \kappa)^{M-1}
\]
where $p_1(\varepsilon, \kappa) = \int_{-\varepsilon/2}^{\varepsilon/2} \frac{e^{\kappa\cos\theta}}{2\pi I_0(\kappa)} d\theta$ is the single-variable window probability under von Mises$(\kappa)$.

\textbf{Remark (alternative reference definitions).} One can also define alignment relative to an anchored reference (e.g., a pacemaker module $\theta_1$) or a random reference. These alternatives yield the same $(M-1)$ scaling exponent, differing only by $O(1)$ constants. We use the mean-phase definition throughout because it matches empirical practice and our simulation protocol.

In all cases:
\begin{equation}\label{eq:coh-mises}
P(A_\varepsilon) \propto p_1(\varepsilon, \kappa)^{M-1}
\end{equation}

\textbf{Small-$\varepsilon$ approximation.} For small windows, the von Mises density near the mode gives $p_1(\varepsilon, \kappa) \approx \varepsilon \cdot e^\kappa / (2\pi I_0(\kappa))$, yielding:
\[
\tau_{\mathrm{coh}} \approx \frac{1}{\lambda_{\mathrm{attempt}}} \left( \frac{2\pi I_0(\kappa)}{\varepsilon \, e^\kappa} \right)^{\alpha(M-1)}
\]
This makes the $\varepsilon^{-(M-1)}$ scaling and the ``exponential in $M$'' dependence explicit. Here $\lambda_{\mathrm{attempt}}$ (s$^{-1}$) is the effective alignment attempt rate, estimated from the decorrelation time of inter-module phase coherence $r_{\mathrm{inter}}(t)$ (see below). \emph{Note:} We use the exact $p_1(\varepsilon, \kappa)$ integral for all numerical estimates (e.g., the binding window calculations use $\varepsilon = \pi/2$); the small-$\varepsilon$ expansion is presented only to clarify the scaling structure.

\textbf{The exponent is $(M-1)$, not $M$.} This arises from quotient geometry: the reference module $\theta_1$ serves as anchor (whether fixed or conditioned upon), leaving $M-1$ independent constraints. This is analogous to the classic result that the probability of $M$ uniform phases falling in \emph{some} semicircle is $M/2^{M-1}$, where the factor $M$ accounts for the free choice of arc location; our anchored definition removes this degree of freedom.

\textbf{Generalization to arbitrary phase constraints.} The alignment event need not be ``all phases near a common reference.'' More generally, suppose a commit requires $k$ independent phase constraints to be satisfied, where each constraint has marginal probability $p_i$ (e.g., specific phase lags between area pairs, cluster relationships, or cross-frequency conditions). Then:
\[
P(\text{alignment}) = \prod_{i=1}^k p_i \approx \bar{p}^k
\]
where $\bar{p}$ is the geometric mean probability. The coherence time scales as $\tau_{\mathrm{coh}} \propto \bar{p}^{-k}$, i.e., exponentially in the number of constraints. Our ``all within $\varepsilon/2$ of mean'' definition is a special case with $k = M-1$ and $p_i = p_1(\varepsilon, \kappa)$. The framework applies whenever commits require multiple semi-independent conditions to be simultaneously satisfied.

\textbf{First-passage time structure.} Assuming independence across modules on the commit timescale, the alignment event $A_\varepsilon$ is a rare target set on the $M$-torus $T^M$. Under fast mixing relative to the target set size, first-passage times are approximately exponential with mean $\tau \approx 1/(\text{attempt rate} \times P(A_\varepsilon))$---a form of Kac's recurrence theorem~\cite{kac1947}.

\begin{proposition}[First-passage time under mixing]\label{prop:fpt}
Let $X(t)$ be an ergodic Markov process on state space $\mathcal{X}$ with stationary distribution $\pi$ and mixing time $\tau_{\mathrm{mix}}$ (the time for $\|P_t(x, \cdot) - \pi\|_{\mathrm{TV}} \leq 1/4$ from worst-case initial state). For a target set $A \subset \mathcal{X}$ with $\pi(A) \ll 1$, the hitting time $T_A = \inf\{t : X(t) \in A\}$ satisfies:
\[
\mathbb{E}[T_A] \approx \frac{\tau_{\mathrm{mix}}}{\pi(A)} \quad \text{as } \pi(A) \to 0.
\]
Moreover, $T_A / \mathbb{E}[T_A]$ is approximately exponentially distributed when $\pi(A) \ll 1$.
\end{proposition}
\noindent This is a standard result in the theory of rare events for Markov chains~\cite{kac1947}. The key assumptions are ergodicity and that the process mixes on timescales much shorter than the expected hitting time. In our setting, ``mixing'' means that successive phase configurations separated by $\tau_{\mathrm{mix}}$ are approximately independent samples from the stationary distribution; hence the Bernoulli-trial heuristic: each mixing interval is an independent ``attempt'' to land in $A_\varepsilon$.

\textbf{Attempt rate.} Let $\tau_{\mathrm{mix}}$ denote the mixing time of the joint phase process---the timescale on which phase configurations become approximately independent of their initial state. The effective number of independent ``alignment attempts'' per unit time is $\sim 1/\tau_{\mathrm{mix}}$. We define the \textbf{attempt rate}:
\[
\lambda_{\mathrm{attempt}} \equiv \frac{1}{\tau_{\mathrm{mix}}} \quad \text{(units: s}^{-1}\text{)}
\]
This is \emph{not} the carrier frequency of oscillations, but the rate at which the relative-phase configuration effectively ``resamples.''

We estimate $\lambda_{\mathrm{attempt}}$ from the inverse decorrelation time of the inter-module order parameter $z_{\mathrm{inter}}(t) = \frac{1}{M}\sum_m e^{i\psi_m(t)}$, where $\psi_m(t)$ is the mean phase of module $m$. The decorrelation time $\tau_{\mathrm{mix}}$ is defined as the $1/e$ crossing of the autocorrelation function of $|z_{\mathrm{inter}}(t)|$, and $\lambda_{\mathrm{attempt}} = 1/\tau_{\mathrm{mix}}$.

\textbf{Simulation proxy.} In simulations, $\lambda_{\mathrm{attempt}}$ is measured directly from the autocorrelation of $|z_{\mathrm{inter}}(t)|$. An alternative proxy---the standard deviation of instantaneous phase velocities across oscillators (``effective frequency spread'' $\Delta\omega$ in rad/s)---is proportional to $\lambda_{\mathrm{attempt}}$ in weakly-coupled regimes, with $\lambda_{\mathrm{attempt}} \approx \Delta\omega$ (since rad is dimensionless). For neural data analysis, we recommend estimating $\lambda_{\mathrm{attempt}}$ from the autocorrelation time of inter-module phase deviations, which directly measures the mixing timescale.

This yields:
\begin{equation}\label{eq:coh-time-exact}
\tau_{\mathrm{coh}} = \frac{1}{\lambda_{\mathrm{attempt}}}\,p_1(\varepsilon, \kappa(r))^{-(M-1)}
\end{equation}

For coupled modules (weak inter-module coupling), the effective exponent is reduced:
\begin{equation}\label{eq:coh-time}
\boxed{\tau_{\mathrm{coh}} \approx \frac{1}{\lambda_{\mathrm{attempt}}}\,p_1(\varepsilon, \kappa(r))^{-\alpha(M-1)}},\qquad \alpha\in(0,1],
\end{equation}
where $\alpha$ captures \emph{effective independence}: $\alpha = 1$ for fully independent modules; $\alpha < 1$ when inter-module coupling induces correlations that reduce the effective number of constraints.

\textbf{Estimating effective constraint dimension.} Define the \textbf{effective constraint dimension}:
\[
d_{\mathrm{eff}} = \alpha(M-1)
\]
This can be estimated from the covariance structure of module phases. Let $C$ be the $M \times M$ covariance matrix of phase deviations $\delta\psi_m = \psi_m - \bar{\psi}$. The \textbf{participation ratio}~\cite{cunningham2014,stringer2019} provides an operational estimate:
\begin{equation}\label{eq:deff-estimator}
\hat{d}_{\mathrm{eff}} = \frac{(\mathrm{tr}\, C)^2}{\mathrm{tr}(C^2)} = \frac{\left(\sum_i \lambda_i\right)^2}{\sum_i \lambda_i^2}
\end{equation}
where $\lambda_i$ are eigenvalues of $C$. This equals $M$ when all eigenvalues are equal (fully independent) and approaches 1 when a single mode dominates (fully correlated). From $\hat{d}_{\mathrm{eff}}$, we obtain $\hat{\alpha} = \hat{d}_{\mathrm{eff}} / (M-1)$.

\textbf{Why correlations affect the exponent.}
For independent phase deviations $\delta\psi_m$ with uniform distribution over $[-\pi, \pi]$,
the probability of all $M$ modules aligning within $\pm\varepsilon/2$ scales as $\sim \varepsilon^{M-1}$.
However, when deviations are correlated with effective rank $d_{\mathrm{eff}}$,
the measure of a small alignment neighborhood scales like $\varepsilon^{d_{\mathrm{eff}}}$
(up to constants depending on the correlation structure).
Hence the exponent is effectively $d_{\mathrm{eff}}$, motivating $\alpha(M-1) \approx d_{\mathrm{eff}}$.
This is analogous to the effective dimensionality of a correlated Gaussian:
if $\mathbf{X} \sim \mathcal{N}(0, \Sigma)$ with eigenvalues $\lambda_i$,
the volume of an $\varepsilon$-ball scales like $\varepsilon^{d_{\mathrm{eff}}}$ where
$d_{\mathrm{eff}} = (\sum_i \lambda_i)^2 / \sum_i \lambda_i^2$ (participation ratio).

\textbf{Log-linear form.} Taking logarithms makes the exponential-in-$M$ scaling explicit:
\[
\log \tau_{\mathrm{coh}} = -\log \lambda_{\mathrm{attempt}} - \alpha(M-1)\log p_1(\varepsilon, \kappa(r)).
\]
Since $p_1 < 1$, we have $-\log p_1 > 0$, so coherence time grows exponentially with effective constraint dimension $\alpha(M-1)$.

\begin{figure}[H]
\centering
\includegraphics[width=\textwidth]{simulations/figures/fig_topology_comparison.pdf}
\caption{\textbf{Validation of coherence time scaling across network topologies.}
(A) Modular, (B) All-to-all, (C) Sparse random networks ($N=100$ oscillators, $K=0.7$, $\varepsilon=1.0$ rad, 20 trials per $M$).
Points show median coherence time $\tau_{\mathrm{coh}}$ with interquartile ranges (error bars).
For heavily censored points (hit fraction $< 0.5$), we report Kaplan-Meier lower bounds with
upward arrows indicating censoring.
We fit Eq.~\ref{eq:coh-time} using two approaches: (i)~\textbf{theory-consistent}:
$\log(\tau \cdot \lambda_{\mathrm{attempt}}) = \alpha(M-1)[-\log p_1(\varepsilon, \kappa(\bar{r}_{\mathrm{inter}}))] + c$
using the trajectory-mean inter-module coherence $\bar{r}_{\mathrm{inter}}$ (averaged over all trials, including censored, to estimate the stationary distribution), and
(ii)~\textbf{linearized surrogate}:
$\log(\tau \cdot \lambda_{\mathrm{attempt}}) \approx \alpha \cdot \tilde{c} \cdot (1-\bar{r}_{\mathrm{inter}})(M-1) + c'$
for comparison.
Dashed lines show fitted predictions with $\hat{\alpha}$ and $R^2$ values indicated.
Modular networks yield $\hat{\alpha} \approx 1$ (near-independent modules, $R^2 = 0.97$); sparse random networks yield $\hat{\alpha} \approx 0.1$ (strong correlations collapse effective constraints); all-to-all coupling yields $\hat{\alpha} \approx 1$ via linearized surrogate (the theory-consistent fit is ill-conditioned in this high-coherence regime where $\bar{r}_{\mathrm{inter}} \approx 1$ makes $-\log p_1 \approx 0$; the surrogate value slightly exceeding 1 reflects linearization error rather than ``more than $M{-}1$ independent constraints'').}
\label{fig:topology}
\end{figure}

\textbf{Key assumptions.} Equation~\ref{eq:coh-time} relies on:
\begin{enumerate}
\item \textbf{Stationarity}: The phase distribution is approximately stationary over the commit timescale.
\item \textbf{Fast mixing}: The joint phase process mixes rapidly relative to the target set size, so first-passage times are approximately exponential (memoryless after conditioning on slowly varying parameters like $r$).
\item \textbf{Weak inter-module coupling}: Modules are semi-independent on the commit timescale; strong coupling would synchronize phases and reduce effective $M$.
\item \textbf{Attempt rate interpretation}: $\lambda_{\mathrm{attempt}}$ represents the inverse correlation time of phase deviations, not the carrier frequency.
\end{enumerate}
\textbf{Diagnostic}: If these assumptions hold, inter-commit intervals (after conditioning on $r$) should be approximately exponentially distributed. Deviations indicate strong coupling, metastability, or non-stationarity. Supplementary simulations confirm this prediction: waiting times conditioned on narrow $r_{\mathrm{inter}}$ bins are approximately exponential (Kolmogorov-Smirnov test), and the von Mises model closure $r \to \kappa \to p_1$ is validated by conditional-on-$r$ fitting ($\kappa_{\mathrm{fitted}}/\kappa_{\mathrm{expected}} = 1.06 \pm 0.12$).

\textbf{Dwell time and the scaling analysis.} Definition~\ref{def:commit} requires alignment to persist for a dwell time $T_{\mathrm{dwell}}$. This introduces a ``hold probability'' $q(T_{\mathrm{dwell}}; r)$---the probability that alignment, once achieved, survives for duration $T_{\mathrm{dwell}}$. The effective commit rate becomes $\lambda_{\mathrm{attempt}} \cdot P(A_\varepsilon) \cdot q$, so $q$ renormalizes the prefactor but does not alter the exponential-in-$M$ scaling, which is driven entirely by $P(A_\varepsilon)$. In simulations we enforce dwell explicitly; for the scaling analysis, $q$ is absorbed into $\lambda_{\mathrm{attempt}}$ as a constant-factor correction and estimated empirically from decorrelation of $|z_{\mathrm{inter}}(t)|$.

\textbf{Simulation limitation.} In the modular Kuramoto validation (Fig.~\ref{fig:topology}), we vary $M$ at fixed total population $N$, which changes module size ($\sim N/M$). This can introduce finite-size effects and alter within- vs between-module coupling balance. For cleaner $M$-scaling, one would keep module size constant by scaling $N$ with $M$. Despite this limitation, the exponential scaling with $M$ is robust across all three topologies tested.

\textbf{Heavy-tailed waiting times.} Empirically observed inter-commit intervals may exhibit heavy tails even when the underlying process is well-described by Eq.~\ref{eq:coh-time}. This arises naturally when $r(t)$, $\lambda_{\mathrm{attempt}}(t)$, or network state drift slowly relative to the commit timescale. The marginal waiting time distribution then becomes a \emph{mixture of exponentials} (a Cox or doubly stochastic process), which can produce power-law-like tails. The key testable prediction is that waiting times \emph{conditional on a narrow bin of $r$} should remain approximately exponential, even when the unconditional distribution is heavy-tailed.

\textbf{Physical interpretation:} $M$ is the \textbf{coordination depth}---the number of semi-independent modules (functional clusters) whose order parameters must jointly exceed a threshold for a commit to register. \emph{Modules are not individual oscillators}: within-module coherence is absorbed into the module's aggregate phase estimate. $\lambda_{\mathrm{attempt}}$ is the \textbf{attempt rate}: the characteristic rate at which inter-module phases decorrelate and ``resample'' (not the carrier frequency). The tolerance $\varepsilon$ is the \textbf{full phase window width} (in radians); $\varepsilon = \pi/2$ corresponds to $\pm 45^\circ$ alignment tolerance, typical for gamma-band binding.

\textbf{Foundational assumption: continuous coordination variables.} This framework focuses on continuous phase variables as the coordination substrate. The commit definition, alignment probabilities, and rare-event statistics are naturally expressed in phase coordinates because phases can \emph{explore} configuration space between commits---the core requirement for the first-passage geometry we model.

Spike trains may participate in coordination (e.g., through spike-field coupling or spike-timing-dependent plasticity), but our framework does not require taking a position on whether spikes or fields are ``fundamental.'' What matters is that \emph{some} continuous coordination variable exists whose alignment determines commit timing. Local field potentials and ephaptic coupling~\cite{anastassiou2011,pinotsis2023} provide one candidate substrate; dendritic integration, subthreshold dynamics, or synaptic state variables could serve similar roles.

\textbf{The skeptical position.} Several arguments suggest ephaptic effects may be computationally negligible:
\begin{enumerate}
\item \emph{Magnitude}: Ephaptic potentials are weak ($\sim$0.1--0.5 mV) compared to synaptic potentials ($\sim$1--10 mV) or action potentials ($\sim$100 mV). The signal is small relative to other influences.
\item \emph{Noise floor}: Neural systems already contend with substantial noise (channel stochasticity, synaptic variability, thermal fluctuations). Ephaptic effects may simply contribute additional noise that gets filtered out.
\item \emph{Redundancy}: Neural circuits have convergent and divergent connectivity. Small timing perturbations in individual neurons should average out across populations.
\item \emph{Synaptic dominance}: Synaptic transmission is the established mechanism of neural communication; ephaptic coupling is at best a second-order correction.
\end{enumerate}

\textbf{Conditions under which ephaptic effects may be significant.} Each skeptical objection above assumes a system operating in a robust, noise-tolerant regime. Near criticality~\cite{beggs2003,shew2013}, the situation may differ:
\begin{enumerate}
\item \emph{Magnitude at criticality}: At a critical point, sensitivity can be enhanced. A 0.1 mV shift may advance or delay spike timing by microseconds; if the system operates near a phase transition, such timing differences could determine which neurons cross threshold first.
\item \emph{Structured vs.\ random noise}: Ephaptic ``noise'' is spatially correlated---neighboring neurons experience similar field perturbations. Redundancy and averaging suppress \emph{uncorrelated} noise; correlated perturbations may propagate coherently.
\item \emph{Timing vs.\ connectivity}: Synaptic weights determine \emph{whether} a neuron eventually fires given sufficient input. Fields influence \emph{when} it fires. For oscillatory coordination, precise timing matters. A neuron that fires 2 ms earlier participates in a different phase relationship.
\item \emph{No central clock}: Biological neural systems lack a global clock; timing emerges from dynamics. The question is not whether fields influence spike timing---they demonstrably do~\cite{anastassiou2011}---but whether those timing shifts propagate to affect function.
\end{enumerate}

We suggest that in unclocked systems operating near criticality, ephaptic coupling \emph{could} be functionally significant. This remains an empirical question. If ephaptic coupling proves to be negligible, analogous coordination constraints would arise via synaptic timing, dendritic integration windows, or other coupling mechanisms---the coherence-time framework applies regardless of the specific coupling substrate. The core claim is that \emph{some} form of inter-module phase coordination limits processing rate; ephaptic coupling is one candidate mechanism, not the only one.

\textbf{The quasi-static coupling regime.} A physical point clarifies the nature of ephaptic coupling. At neural oscillation frequencies ($\sim$1--100 Hz), electromagnetic fields in brain tissue are quasi-static: even accounting for the enormous relative permittivity of neural tissue ($\varepsilon_r \sim 4.5 \times 10^7$ at 40 Hz~\cite{gabriel1996}), the electromagnetic phase velocity remains $\sim$44 km/s, giving wavelengths of order 1 km---far larger than any neural structure. Ephaptic coupling therefore operates in the near-field regime where retardation effects are negligible; field changes are felt essentially instantaneously across cortical columns.

The empirically observed propagation speed of ephaptic effects ($\sim$0.1 m/s through transected tissue where synaptic transmission is excluded~\cite{subramanian2022}) is not an electromagnetic wave velocity but a \emph{synchronization front}: each oscillator must be dragged past a threshold by its neighbours' fields before it can influence the next, producing a relay-like cascade much slower than the underlying field propagation. This distinction matters for the coherence time framework because it means inter-oscillator coupling is effectively instantaneous (quasi-static fields), while the \emph{spread of coherence} through a population is limited by oscillator dynamics, not by field propagation delays.

\textbf{Regional variation in ephaptic coupling.} Notably, ephaptic coupling operates differently across brain regions. Cerebellar architecture, with its densely packed parallel fibers and Purkinje cells, supports strong local field effects~\cite{han2018}. Cortical architecture, with its columnar organization and heterogeneous cell types, produces different field geometries. This regional variation in ephaptic coupling strength may explain why cerebellar processing appears exempt from the coordination penalties we attribute to high-$M$ cortical assemblies---the cerebellum may use different coupling mechanisms that bypass the phase-alignment bottleneck.

\subsection{The speed-flexibility trade-off}

The coherence time reveals opposing constraints on commit rate:

\begin{proposition}[Speed-Flexibility Frontier]\label{prop:tradeoff}
At fixed power $P$, systems face a Pareto frontier in $(M,r)$ space:
\begin{itemize}
\item \textbf{Increasing $M$}: Expands combinatorial flexibility (number of modules whose phases can vary independently) but slows commits via $\tau_{\mathrm{coh}} \propto (\mathrm{const})^{M}$
\item \textbf{Increasing $r$}: Speeds commits via reduced circular variance $(1-r)$ but restricts exploration to low-dimensional synchronized manifold
\end{itemize}
\end{proposition}

Increasing coordination depth $M$ expands the dimensionality of commit-registered outcomes but requires exponentially longer alignment time. Increasing $r_{\mathrm{inter}}$ reduces the dimensionality of the \textbf{relative-phase degrees of freedom}, constraining dynamics toward a lower-dimensional \textbf{phase-synchronized manifold}~\cite{strogatz2000,ott2008}. This makes multi-module coordination faster (larger $p_1$) but reduces the number of effectively independent constraints (smaller effective $\alpha$).

Importantly, this does not imply that \textbf{all} neural activity becomes low-dimensional; amplitude patterns, cross-frequency interactions, and within-module heterogeneity can remain rich. The dimensionality reduction refers specifically to the inter-module phase coordination subspace.

\subsection{Quantum and thermodynamic limits of computation}

A natural question is how this coordination-based framework relates to fundamental physical limits on computation. We now explicitly address the quantum and thermodynamic bounds that form the basis of information transmission in living systems.

\subsubsection{The Landauer bound}

Landauer's principle~\cite{landauer1961} establishes that any logically irreversible operation---specifically, the erasure or overwriting of one bit of information---dissipates at least $\kB T \ln 2$ of energy as heat. This bound arises from the connection between logical irreversibility and thermodynamic entropy: erasing a bit reduces the system's information-theoretic entropy, which must be compensated by entropy increase in the environment~\cite{bormashenko2024}.

The minimum energy per bit at biological temperature ($T \approx 310$ K) is:
\[
E_{\mathrm{Landauer}} = \kB T \ln 2 \approx 3 \times 10^{-21} \text{ J/bit}
\]

This sets an absolute floor on energy dissipation per irreversible operation. However, the Landauer bound constrains energy per operation, not time per operation. A system could in principle perform operations arbitrarily fast if sufficient power were available---until quantum speed limits intervene.

\subsubsection{The Margolus-Levitin bound and quantum speed limits}

The Margolus-Levitin theorem~\cite{margolus1998} establishes that the minimum time to transition between orthogonal quantum states is:
\[
\tau_{\mathrm{ML}} = \frac{\pi\hbar}{2\langle E \rangle}
\]
where $\langle E \rangle$ is the average energy above the ground state. This bound, together with the Mandelstam-Tamm relation~\cite{mandelstam1945} ($\tau_{\mathrm{MT}} = \hbar/(2\Delta E)$ where $\Delta E$ is energy uncertainty), constitutes the quantum speed limit on state evolution.

These bounds originate from the energy-time uncertainty relation $\Delta E \cdot \Delta t \gtrsim \hbar/2$, which has been extensively analyzed in the context of biological computation and quantum-like modeling~\cite{khrennikov2021,igamberdiev2025constraints,igamberdiev2026autopoietic}. As Bormashenko~\cite{bormashenko2024} notes, synthesis of the Landauer bound with the Margolus-Levitin limit yields a minimal computation time scaling as $\tau_{\min} \sim h/\kB T$---decreasing temperature reduces energy consumption per operation but simultaneously slows computation.

At biological temperatures with $\langle E \rangle \sim \kB T \approx 4 \times 10^{-21}$ J:
\[
\tau_{\mathrm{QSL}} \sim \frac{\pi \times 10^{-34}}{2 \times 4 \times 10^{-21}} \sim 10^{-13} \text{ s}
\]

This is fourteen orders of magnitude faster than the millisecond timescales of neural processing.

\subsubsection{Why coherence time dominates in biological systems}

The central claim of this paper is that in biological neural systems, the quantum speed limit and Landauer bound are typically \textbf{non-binding constraints}. The dominant bottleneck is coherence time---the time required for distributed oscillatory assemblies to achieve sufficient phase coordination for a commit to occur.

Consider the hierarchy of timescales for cortical visual processing:
\begin{align}
\tau_{\mathrm{QSL}} &\sim 10^{-13} \text{ s} \quad \text{(quantum speed limit)} \\
\tau_{\mathrm{SNR}} &\sim 10^{-3} \text{ s} \quad \text{(signal detection)} \\
\tau_{\mathrm{coh}} &\sim 10^{-2} \text{ s} \quad \text{(coordination time)}
\end{align}
(Note: The Landauer bound constrains \emph{energy} per bit, not time; we omit it from this hierarchy since the time-limiting step is coordination, not dissipation.)

The coherence time exceeds quantum limits by eleven orders of magnitude. This enormous gap arises because neural computation is not limited by how fast individual elements can switch states, but by how long it takes for \emph{distributed assemblies of semi-independent modules} to achieve coordinated phase alignment.

Mathematically, this reflects the exponential scaling of coherence time with coordination depth:
\[
\tau_{\mathrm{coh}} \propto p_1^{-\alpha(M-1)}
\]
For $M = 8$ modules with $p_1 \approx 0.7$ and $\alpha = 0.7$, the exponential factor is $(0.7)^{-4.9} \approx 6$. Even modest coordination requirements ($M \sim 5$--15) produce millisecond-scale coherence times from microsecond-scale phase dynamics.

This analysis clarifies the relationship between our framework and fundamental physics: quantum and thermodynamic limits set absolute bounds on computation, but biological systems operate in a regime where multi-module coordination---not quantum switching speed or Landauer dissipation---determines processing rate.

\subsection{Parameter estimation protocols}

\paragraph{Coordination depth ($M$).} $M$ is the count of semi-independent modules whose order parameters must jointly exceed a threshold to register a commit.
\textbf{Operational estimation:}
\begin{enumerate}
\item Cluster channels/units by band-limited coherence or Granger causality into modules.
\item Identify the minimal set of modules whose joint activity predicts behavioral commits.
\item Typical values: occipital visual tasks yield $M \sim 6$--10; cross-modal binding $M \sim 10$--15.
\end{enumerate}

\paragraph{Kuramoto coherence ($r$).} Extract instantaneous phase $\theta_j(t)$ from bandpass-filtered LFP/EEG via Hilbert transform. For the coherence time model, we distinguish:
\begin{itemize}
\item \textbf{Within-module coherence:} $r_{\mathrm{intra},m}(t) = \left|\frac{1}{n_m}\sum_{j \in m} e^{i\theta_j(t)}\right|$ for module $m$ with $n_m$ channels.
\item \textbf{Inter-module coherence:} $r_{\mathrm{inter}}(t) = \left|\frac{1}{M}\sum_{m=1}^M e^{i\psi_m(t)}\right|$ where $\psi_m$ is module $m$'s aggregate phase.
\end{itemize}
Commits require high $r_{\mathrm{intra}}$ (modules are internally coherent) \emph{and} rare alignment of inter-module phases (low baseline $r_{\mathrm{inter}}$, with commits occurring when $r_{\mathrm{inter}}$ transiently exceeds threshold). The $r$ and $\kappa$ in Eq.~\ref{eq:coh-time} refer to the \emph{inter-module} distribution. Typical inter-module values: spontaneous cortex $\sim 0.2$--0.4; attention $\sim 0.5$--0.7~\cite{fries2015,womelsdorf2007}.

\textbf{Interpretation of $r$ in Eq.~\ref{eq:coh-time}.} The parameter $r$ can be interpreted in two ways: (1) as a slowly varying \emph{state variable} $r(t)$, in which case Eq.~\ref{eq:coh-time} gives an instantaneous hazard rate $h(t) = \lambda_{\mathrm{attempt}}\, p_1^{\alpha(M-1)}$, and the waiting time distribution becomes a hazard integral over the trajectory of $r(t)$; or (2) as a \emph{conditional} parameter representing coherence during the pre-commit interval, in which case Eq.~\ref{eq:coh-time} gives the mean waiting time conditional on that coherence level. Interpretation (1) naturally explains heavy-tailed waiting time distributions when $r(t)$ drifts slowly.

\paragraph{Phase alignment window ($\varepsilon$).} From spike train or LFP cross-correlograms, measure half-width at half-maximum (HWHM, units: seconds). The \textbf{full} phase window is $\varepsilon = 2\,\omega_0\,\mathrm{HWHM}$ (radians).

\paragraph{Attempt rate ($\lambda_{\mathrm{attempt}}$).} Estimate from the decorrelation time of inter-module coherence $|z_{\mathrm{inter}}(t)|$: $\lambda_{\mathrm{attempt}} = 1/\tau_{\mathrm{mix}}$. Alternatively, estimate from instantaneous frequency spread $\Delta\omega$ (rad/s): since radians are dimensionless, $\lambda_{\mathrm{attempt}} \approx \Delta\omega$ s$^{-1}$. Typical values for cortical alpha/beta bands: $\lambda_{\mathrm{attempt}} \sim 60$--$200$ s$^{-1}$ (corresponding to frequency spreads of $\sim 10$--$30$ Hz).

\textbf{Analysis workflow summary.} To apply the coherence time framework to neural data:
\begin{enumerate}
\item \textbf{Define modules}: Use community detection on coherence matrices, or anatomical ROIs.
\item \textbf{Compute module phases}: Extract $\psi_m(t)$ from bandpass-filtered signals via Hilbert transform; verify within-module coherence $r_{\mathrm{intra},m} \geq r_{\mathrm{thr}}$.
\item \textbf{Compute inter-module coherence}: $r_{\mathrm{inter}}(t) = |M^{-1}\sum_m e^{i\psi_m(t)}|$; convert to $\kappa$ via $r = I_1(\kappa)/I_0(\kappa)$.
\item \textbf{Estimate $\varepsilon$}: Measure cross-correlogram HWHM; $\varepsilon = 2\omega_0 \cdot \mathrm{HWHM}$.
\item \textbf{Estimate $\lambda_{\mathrm{attempt}}$}: Compute autocorrelation of $|z_{\mathrm{inter}}(t)|$; set $\lambda_{\mathrm{attempt}} = 1/\tau_{\mathrm{mix}}$ where $\tau_{\mathrm{mix}}$ is the $1/e$ crossing.
\item \textbf{Fit $\alpha$}: Either from participation ratio of module-phase covariance (Eq.~\ref{eq:deff-estimator}), or from scaling of $\tau$ with $M$ across task conditions.
\end{enumerate}

\section{Results}

\subsection{Visual perceptual binding windows}

\textbf{Censoring note.} For simulation trials that did not reach the alignment threshold within simulation time $T$, we treat these as right-censored observations and report Kaplan-Meier survival estimates of $\tau_{\mathrm{coh}}$. Points with hit fraction $< 0.5$ are shown with upward arrows indicating the median is censored beyond $T$.

Human visual perception exhibits temporal integration windows of 30--50 ms for perceptual binding tasks. Note that critical flicker fusion (CFF) thresholds are higher ($\sim$60 Hz photopic) for simple luminance detection---a low-$M$ task requiring minimal inter-module coordination. The gap between CFF and binding-window timescales is itself predicted by the framework: low-$M$ commits are fast; high-$M$ perceptual binding is slow. We apply Eq.~\ref{eq:main-bound} with neural parameters for the high-$M$ binding case.

\textbf{Parameter provenance:} We take $M=8$ occipital modules (a plausible count of functional submodules in visual binding, including V1, V2, V4, and posterior parietal subdivisions; can be estimated via community detection on coherence networks). Coherence $r=0.7$--$0.8$ reflects attentional engagement~\cite{fries2015}. Full tolerance $\varepsilon=\pi/2$ rad ($\pm 45^\circ$) corresponds to the phase windows observed in gamma-band cross-correlograms during visual binding; note that $\pm 90^\circ$ tolerance would be too loose for functional synchrony. Attempt rate $\lambda_{\mathrm{attempt}} \approx 126$ s$^{-1}$ is estimated from the decorrelation time of inter-module coherence in alpha/beta bands (corresponding to a frequency spread of $\sim 20$ Hz; not the carrier frequency). Effective independence $\alpha=0.7$ reflects modular cortical topology; it can be fit empirically from the scaling of inter-commit times with $M$.

First compute $p_1$: at $r=0.7$, $\kappa \approx 2.0$, giving $p_1(\pi/2, 2.0) \approx 0.68$ (numerical integration of Eq.~\ref{eq:coh-mises}). Using Eq.~\ref{eq:coh-time} with $\lambda_{\mathrm{attempt}} \approx 126$ s$^{-1}$ and $\alpha=0.7$:
\[
\tau_{\mathrm{coh}} \approx \frac{1}{126}\,(0.68)^{-0.7 \times 7}
= \frac{1}{126}\,(0.68)^{-4.9} \approx 54~\mathrm{ms}.
\]
With higher coherence $r=0.8$ ($\kappa \approx 2.9$, $p_1(\pi/2, 2.9) \approx 0.78$):
\[
\tau_{\mathrm{coh}} \approx \frac{1}{126}\,(0.78)^{-4.9} \approx 27~\mathrm{ms}.
\]

Computing all terms in Eq.~\ref{eq:main-bound}:
\begin{align}
\tau_{\mathrm{QSL}} &\sim 10^{-13} \,\mathrm{s} \quad \text{(negligible---quantum limit non-binding)} \\
\tau_{\mathrm{SNR}} &\sim 5\text{--}10~\mathrm{ms} \quad \text{(task-dependent)} \\
\tau_{\mathrm{power}} &\ll 1~\mathrm{ms} \quad \text{(per commit)} \\
\tau_{\mathrm{coh}} &\approx 27\text{--}54~\mathrm{ms} \quad \text{($r = 0.7$--$0.8$, $M = 8$)}
\end{align}

The effective commit time is $\tau_{\mathrm{eff}} \gtrsim \max[\text{QSL, SNR, power, coh}] \approx 27$--$54~\mathrm{ms}$, consistent with human binding windows (30--50 ms) given empirically plausible parameter ranges. \textbf{Coherence time dominates}; quantum and Landauer bounds are not the limiting factors. This is an order-of-magnitude existence proof---the framework yields the correct timescale with independently constrained parameters---not a precise prediction of a single value.

\paragraph{Sensitivity analysis.} Table~\ref{tab:sensitivity} shows how $\tau_{\mathrm{coh}}$ varies across plausible parameter ranges. The 20--80 ms binding window is robust to parameter uncertainty. All $p_1$ values computed via numerical integration of Eq.~\ref{eq:coh-mises}; $\kappa$ obtained from $r$ via numerical inversion of $r = I_1(\kappa)/I_0(\kappa)$.

\textbf{Sanity check:} For any $\kappa \geq 0$, $p_1(\pi, \kappa) \geq 0.5$ with strict inequality for $\kappa > 0$, since the window $\varepsilon = \pi$ covers half the circle centered on the mode of a unimodal distribution.

\begin{table}[H]
\centering
\caption{Sensitivity of coherence time to parameter variations. Base case: $M=8$, $r=0.75$, $\varepsilon=\pi/2$, $\lambda_{\mathrm{attempt}} = 126$ s$^{-1}$, $\alpha=0.7$.}
\label{tab:sensitivity}
\begin{tabular}{lccccc}
\toprule
Parameter varied & $M$ & $r$ & $\varepsilon$ & $\lambda_{\mathrm{attempt}}$ (s$^{-1}$) & $\tau_{\mathrm{coh}}$ (ms) \\
\midrule
Base case & 8 & 0.75 & $\pi/2$ & 126 & 39 \\
Lower $M$ & 6 & 0.75 & $\pi/2$ & 126 & 25 \\
Higher $M$ & 10 & 0.75 & $\pi/2$ & 126 & 61 \\
Lower $r$ & 8 & 0.65 & $\pi/2$ & 126 & 75 \\
Higher $r$ & 8 & 0.85 & $\pi/2$ & 126 & 19 \\
Narrower $\varepsilon$ & 8 & 0.75 & $\pi/3$ & 126 & 105 \\
Faster $\lambda_{\mathrm{attempt}}$ & 8 & 0.75 & $\pi/2$ & 189 & 26 \\
\bottomrule
\end{tabular}
\end{table}

The exponential dependence on $M$ and strong dependence on $r$ mean that modest changes in attention state (affecting $r$) or task complexity (affecting $M$) can shift binding windows by factors of 2--4, consistent with empirical observations of attention-dependent temporal resolution~\cite{samaha2015}.

\subsection{Tachypsychia: a dual-loop hypothesis}

During acute stress or falls, subjects report subjective time slowing while objective temporal resolution remains unchanged~\cite{stetson2007}. We propose a \textbf{dual-commit architecture} as one explanation for this dissociation (alternative accounts exist, including attentional sampling rate changes and memory density effects):

\emph{Perceptual loop} (cortical): High-$M$ ($\sim 10$--15 modules) coherent field dynamics. Commits sparse (5--20 Hz). Arousal increases coherence $r$ and thus information rate per commit interval.

\emph{Motor loop} (cerebellar/basal ganglia): Low-$M$ ($\sim 3$--5 modules) primitives. Motor latencies of 50--150 ms reflect simpler coordination requirements (lower $M$, higher $\alpha$) than multi-area perceptual binding. Arousal modulates decision threshold, preserving reaction time.

\textbf{Measurable proxies:} Perceptual loop commit rate could be indexed by temporal-order judgment thresholds or EEG phase coherence in theta/alpha bands. Motor loop could be indexed by simple RT distributions or EMG onset latencies. The hypothesis predicts these measures should be \emph{statistically dissociable} under arousal manipulation, with temporal-order thresholds changing while simple RT remains stable.

\subsection{Alpha oscillations and temporal acuity}

Alpha oscillations (8--12 Hz) in visual cortex correlate with temporal acuity: individuals with faster alpha perceive time faster~\cite{samaha2015,cecere2015}. One interpretation is that alpha acts as a gating rhythm that modulates commit opportunity windows, effectively setting the attempt rate $\lambda_{\mathrm{attempt}}$ or the sampling schedule for phase alignment.

\textbf{Prediction:} If alpha frequency modulates commit timing, then entraining alpha via transcranial alternating current stimulation (tACS)~\cite{helfrich2014} should shift perceptual binding windows approximately linearly with frequency change. This prediction distinguishes the gating interpretation from alternatives where alpha reflects post-commit memory encoding.

\subsection{Metabolic scaling across species}

When power is limiting ($\tau_{\mathrm{power}} \gg \tau_{\mathrm{coh}}$), we obtain $\tau_{\mathrm{eff}} \sim 1/P$. Critical flicker fusion frequency should scale with mass-specific metabolic rate.

Critical flicker fusion frequency varies enormously across taxa: flies $\sim 240$ Hz, humans $\sim 60$ Hz, leatherback turtles $\sim 15$ Hz. Healy et al.~\cite{healy2013} found that mass-specific metabolic rate explains substantial variance in CFF across 34 vertebrate species spanning three orders of magnitude in body mass, consistent with the predicted power-limited regime for small, high-metabolism animals.

\section{Discussion}

\subsection{Biological examples of the speed-flexibility trade-off}

\paragraph{Mind-wandering and creative insight (high $M$, low $r$).}
Default-mode network activity~\cite{raichle2001} exhibits low inter-regional coherence ($r \sim 0.3$) and high coordination depth ($M \sim 12$--20). Commits are rare ($\sim 0.1$--1 Hz), producing subjective ``slow thinking'' but enabling novel associations.

\paragraph{Focused attention (moderate $M$, moderate $r$).}
In our model, attention increases coherence ($r: 0.3 \to 0.6$) and narrows coordination to task-relevant modules ($M: 15 \to 8$) (model parameters motivated by the communication-through-coherence framework~\cite{fries2015}). This produces faster commits (10--30 Hz) in a restricted subspace.

\paragraph{Automaticity and skill learning (low $M$, high $r$).}
Overlearned motor sequences compress to low-dimensional primitives ($M \sim 3$--5) with tight coherence ($r \sim 0.8$)~\cite{wolpert1998}. Commits are rapid (50--150 ms) but inflexible.

\subsection{Temporal embedding as a pre-commit representation mechanism (hypothesis)}

\textbf{Why include this section.}
The coherence time bound has implications beyond the derivation of Eq.~\ref{eq:coh-time}:
it constrains the useful window over which \textbf{temporally embedded representations}
remain reconstructible. If pre-commit intervals exceed $\tau_{\mathrm{coh}}$,
downstream readout mechanisms cannot reliably infer earlier states.
Cerebellar delay lines are a plausible biological implementation of preserving
distributed state across these pre-commit intervals. This section illustrates
one potential design strategy for coping with finite coherence times, but
\textbf{is not required for the derivation or validation of Eq.~\ref{eq:coh-time}}.
It is included to demonstrate practical consequences of the bound for
neural system design.

Coherence time tells us \emph{how long} pre-commit dynamics must be maintained before alignment occurs. A separate question is \emph{how} high-dimensional structure can remain accessible during that interval. We propose---as a hypothesis for future investigation---that biological systems may implement \textbf{Takens delay embedding}~\cite{takens1981}.

\paragraph{Takens' theorem.}
For a deterministic system with state dimension $d$, time-delayed copies of a scalar observable---$x(t), x(t-\tau), x(t-2\tau), \ldots$---can reconstruct the complete state-space topology. Crucially, $2d+1$ delay coordinates generically suffice.

\paragraph{A candidate biological implementation: cerebellar parallel fibers.}
Parallel fibers---unmyelinated granule cell axons---conduct at 0.2--0.5 m/s with systematically varying path lengths, creating conduction delays of 5--15 ms~\cite{braitenberg1958,llinas1982,wyatt2005}. These delays match the optimal embedding parameter $\tau^* = 1/(4f)$ for motor dynamics at 15--50 Hz.

We hypothesize that the cerebellum may provide one biological implementation of delay-embedded state preservation during pre-commit intervals. This would enable rapid low-dimensional readout (motor commands) from high-dimensional state without requiring simultaneous global coordination. \textbf{This remains conjecture}; direct evidence would require demonstrating that cerebellar computations exploit delay structure for state reconstruction.

\paragraph{Robustness to noise.}
We use Takens embedding as a \emph{design principle} rather than a literal guarantee: delay-line architectures improve state inference under partial observability, even when exact reconstruction is impossible.

\paragraph{Connection to coherence time.}
Coherence time determines the maximum window over which temporal embeddings remain usable. If coherence decays faster than the embedding window, reconstruction fails and high-dimensional structure is lost before a commit can occur.

\begin{figure}[H]
\centering
\includegraphics[width=\textwidth]{figures/cerebellar_takens.pdf}
\caption{\textbf{Delay embedding reconstructs high-dimensional dynamics from scalar observation.} Simulation demonstrating Takens' theorem applied to cerebellar architecture. (a-b) Lorenz attractor (ground truth, $d \approx 2.06$) and its 1D projection (what a slow cortical oscillation might carry). (c-d) Delay embedding with $\tau=8$, $m=3$ reconstructs attractor topology from the scalar signal alone. (e) Neighborhood preservation quantifies reconstruction quality across embedding dimensions. (f) Theoretical implications: parallel fibers with graded conduction delays implement delay embedding, enabling reconstruction of high-dimensional coordination signals from low-frequency input.}
\label{fig:takens}
\end{figure}

\subsection{Generalization beyond neural systems}

While we emphasize neural applications given their empirical accessibility, the coherence time framework applies to any biological system whose function requires coordinated phase relations and whose outputs are sparse commitment events.

\paragraph{Scope statement.} Neural systems are the best-instrumented instance of a broader class: living, self-maintaining oscillator assemblies that perform irreversible state registration. We do not treat collectives (e.g., societies) as intelligent systems in this sense; they are coordination networks of intelligent organisms and introduce additional channels not modeled here.

\paragraph{Chromosomal coordination.} Chromosomal systems provide a non-neural example of the same architecture. High-dimensional continuous dynamics (polymer conformations, electrostatic interactions, topological constraints, phase separation) produce function via sparse low-dimensional commits (origin firing, checkpoint decisions, segregation outcomes).

Anaphase onset is naturally modeled as a commit event: a global irreversible transition triggered only after multiple semi-independent units satisfy a coherence-like constraint. The waiting time should scale steeply with effective coordination depth, consistent with first-passage into a small target set. The relevant ``order parameter'' is any coordination metric determining whether a global event can proceed---spatial colocalization, topological state, mechanical alignment---not necessarily Kuramoto $r$.

\paragraph{Circadian networks.} Circadian gene networks exhibit $M \sim 3$--5 core feedback loops, $r \sim 0.8$ (tight coupling), and $\varepsilon \sim 0.5$ rad (loose phase tolerance). These parameters yield coordination times much shorter than the 24-hour period, implying coordination is \emph{not} the bottleneck. The 24-hour period arises from molecular lags. The framework correctly identifies the dominant regime: circadian clocks are biochemical-delay-limited, not coordination-limited.

\subsection{Testable predictions}

\paragraph{Alpha entrainment.}
tACS at alpha frequency should shift perceptual binding windows linearly with frequency change. 20\% frequency increase predicts 20\% threshold reduction.

\paragraph{Arousal effects on dual-loop dissociation.}
During simultaneous simple reaction time + temporal-order judgment under acute stress:
\begin{itemize}
\item Temporal-order discrimination threshold changes
\item Simple reaction time unchanged or slightly faster
\item Statistical independence between measures
\end{itemize}

\paragraph{Coherence measures and consciousness.}
If phenomenal experience reflects pre-commit coherent dynamics, conscious states should correlate more strongly with phase-locking index and theta-gamma coupling than with spike count or mean firing rate.

\subsection{Relationship to existing frameworks}

\textbf{Quantum-like cognition.} Recent work has explored connections between oscillatory neural networks and quantum-like cognitive phenomena~\cite{khrennikov2025bridge}. Our framework provides a classical high-dimensional explanation for why such quantum-like signatures emerge: high-dimensional dynamics with limited measurement bandwidth produce phenomenology parallel to quantum systems---measurement-like collapse at commits, superposition-like coexistence of states pre-commit, uncertainty-like timing inaccessibility.

\textbf{Integrated Information Theory}~\cite{tononi2016}: We provide thermodynamic/temporal foundations. Integration costs energy and takes time (exponentially with coordination depth).

\textbf{Free Energy Principle}~\cite{friston2010}: We add temporal structure. Prediction updates occur at commit events with rate limited by Eq.~\ref{eq:main-bound}.

\subsection{Frequency and dimensionality: two interpretations}

A prominent view in systems neuroscience holds that low-frequency oscillations (delta, theta, alpha) support low-dimensional, ``simple'' dynamics, while high-frequency oscillations (gamma) support high-dimensional, ``complex'' computations~\cite{miller2018,lundqvist2016}. Miller and colleagues argue that slow oscillations provide a low-dimensional ``scaffold'' or ``broadcast'' signal, while fast gamma carries detailed content in a higher-dimensional space. Recent work extends this framework to cytoelectric coupling, proposing that electric fields organize neural activity via ephaptic effects down to the cytoskeletal level~\cite{pinotsis2023}.

\textbf{An alternative interpretation.} The coherence time framework suggests a different relationship between frequency and dimensionality, based on coordination constraints:

\paragraph{The physics argument.} Consider an oscillatory system that must achieve phase alignment to register a commit. The number of alignment ``attempts'' per oscillation cycle is roughly $\lambda_{\mathrm{attempt}} / f$, where $f$ is the carrier frequency and $\lambda_{\mathrm{attempt}}$ is the attempt rate. For alignment probability $p_1$:
\begin{itemize}
\item \textbf{High frequency}: Few attempts per cycle. To achieve reliable alignment within one period, the system requires high $p_1$, which requires high coherence $r$. But high $r$ confines dynamics to a low-dimensional synchronized manifold. \emph{Fast oscillations force low dimensionality.}
\item \textbf{Low frequency}: Many attempts per cycle. The system can tolerate low $p_1$ and thus low $r$. Phases can explore more of configuration space while still achieving alignment before the next cycle. \emph{Slow oscillations permit high dimensionality.}
\end{itemize}

This is not a correlation but a constraint: maintaining high-dimensional coherent dynamics at high frequency is physically harder because there is less time for the rare alignment events that preserve high-$M$ coordination.

\paragraph{The substrate mechanism: dimensionality in phase offsets.} The argument above can be made more precise. In an assembly of coupled oscillators, representational capacity resides in the \emph{phase offsets} between oscillators---the pattern of relative timing that constitutes the system's high-dimensional state. The number of distinguishable phase configurations scales as $d_{\mathrm{eff}} \sim T / \tau_{\mathrm{noise}}$, where $T = 1/f$ is the oscillation period and $\tau_{\mathrm{noise}}$ is the temporal precision set by the substrate noise floor. At lower frequency, each cycle provides a longer interval over which phase offsets can be resolved, permitting more distinguishable configurations and thus higher effective dimensionality. This gives $d_{\mathrm{eff}} \sim 1/(f \cdot \tau_{\mathrm{noise}})$: dimensionality is inversely proportional to frequency, with the noise floor as the key parameter.

This relationship has implications beyond neuroscience. For any physical oscillator assembly---biological or engineered---the representational capacity per cycle is bounded by the ratio of period to timing noise. Engineering lower noise floors (e.g., in neuromorphic substrates) increases dimensionality at any given frequency, while operating at lower frequencies achieves the same effect without substrate improvement. The biological nervous system appears to exploit both strategies: low-frequency cortical oscillations for high-dimensional integration, and the cerebellum's low-noise, tightly coupled architecture for rapid low-dimensional processing.

\paragraph{Reinterpreting the empirical evidence.} The observation that gamma activity correlates with ``detailed'' processing and alpha/theta with ``global'' states is consistent with our framework but admits a different interpretation. High-frequency gamma may reflect \emph{local}, low-$M$ modules committing rapidly---not high-dimensional integration but many parallel low-dimensional commits. Meanwhile, low-frequency oscillations may coordinate high-$M$ assemblies that \emph{require} slow timescales to achieve alignment.

\paragraph{The observability problem.} Claims about neural dimensionality face a fundamental measurement limitation: electrodes are sparse, noisy, low-dimensional sensors. One can only estimate dimensionality up to the rank of the observation matrix. If slow oscillations genuinely support high-dimensional dynamics (as our framework predicts), this would not necessarily be visible through electrode recordings---the measurement apparatus itself constrains what can be observed.

Consider an analogy: the gravitational field at a point is a 3-dimensional vector, yet it integrates information about the distribution of all mass in the universe. The low-dimensional measurement (a vector) does not imply low-dimensional underlying structure. Similarly, a slow oscillation measured as a scalar or low-rank signal may nonetheless coordinate high-dimensional dynamics that project onto that measurement. The claim that ``gamma is high-dimensional'' may reduce to ``gamma is the highest-dimensional signal we can resolve with our electrodes''---a statement about the measurement, not the system.

This reflects a general methodological point: physics uses low-dimensional mathematics because it is computable and because our instruments measure low-dimensional projections. The predictive success of low-dimensional models demonstrates that low-dimensional descriptions are \emph{useful}, not that nature is low-dimensional. The same applies to neuroscience: the success of low-rank decompositions (PCA, factor analysis, demixed PCA) in capturing behaviorally relevant variance does not establish that neural dynamics are intrinsically low-dimensional---it establishes that low-dimensional projections correlate with the behavioral variables we measure. These are epistemological claims about our models, not ontological claims about the system.

\paragraph{Testable predictions distinguishing the two views.}
\begin{enumerate}
\item \textbf{Dimensionality-frequency relationship}: If Miller et al.\ are correct, pharmacologically or optogenetically increasing gamma power should increase state-space dimensionality (e.g., participation ratio of neural covariance). If our framework is correct, increasing gamma power while holding coordination depth $M$ constant should \emph{decrease} dimensionality (forcing tighter coupling to maintain coordination at higher frequency).
\item \textbf{Information content per cycle}: Miller et al.\ predict more bits per gamma cycle than per alpha cycle. We predict the opposite: more bits per alpha cycle (because high-$M$ commits are possible), fewer bits per gamma cycle (because only low-$M$ commits can keep pace).
\item \textbf{Cross-frequency coupling direction}: If slow oscillations scaffold fast ones (Miller), gamma amplitude should be phase-locked to theta/alpha with high fidelity. If fast oscillations are dimensionally constrained (our view), gamma bursts should occur when local modules achieve rapid low-$M$ commits, potentially \emph{out of phase} with slow oscillations that coordinate broader assemblies.
\end{enumerate}

We note that this disagreement is empirically tractable. The coherence time framework makes quantitative predictions about the relationship between oscillation frequency, coupling strength, and state-space dimensionality that can be tested with existing recording and analysis methods.

\paragraph{Requisite variety and nonlocal dimensionality constraints.} The frequency-dimensionality relationship is further complicated by a control-theoretic constraint. Ashby's law of requisite variety~\cite{ashby1956} states that a controller must have at least as much variety (effective dimensionality) as the system it regulates. In a coupled oscillator assembly, this means that a subnetwork operating in a low-dimensional regime---whether due to high coherence, pathological constraint, or deliberate dimensional compression---loses the degrees of freedom needed to regulate the regions it is coupled to. The dimensions of the controlled system that are orthogonal to the controller's accessible subspace become effectively unregulated.

This has two consequences for interpreting neural dimensionality. First, dimensionality is not a purely local property: constraining one region to a low-dimensional attractor has nonlocal effects on the controllability of coupled regions, potentially producing high-dimensional wandering in areas that are not themselves pathological but have lost their regulatory input. Second, simple claims about ``high-dimensional'' or ``low-dimensional'' processing in a given brain region must account for the coupling topology---the same region may exhibit different effective dimensionality depending on the state of its controllers, not only its own intrinsic dynamics. The speed-flexibility trade-off derived in \S\ref{prop:tradeoff} is, in this light, a temporal expression of Ashby's constraint: high-$M$ coordination provides the requisite variety for flexible regulation but pays an exponential cost in commit time.

\subsection{Implications for artificial systems}

The coherence time framework identifies a fundamental design constraint for any distributed computing system that lacks a global clock.

\paragraph{Clocked architectures.} Conventional digital computers (CPUs, GPUs, current AI accelerators) enforce synchronization via a global clock signal. Every operation is gated by a central timing reference; state updates occur in lockstep across all components. This architecture sidesteps the coherence time penalty entirely---there is no waiting for spontaneous phase alignment because alignment is \emph{imposed} by the clock.

We hypothesize that this comes at a cost beyond power and wire delays: \textbf{clocked systems may be constrained in ontological dimensionality}. The clock imposes a global update constraint that can suppress the kinds of metastable, spontaneously coordinated dynamics we model here. Even massively parallel operations (matrix multiplications across thousands of cores) proceed via synchronized state updates---the system can \emph{represent} high-dimensional states in memory, but may not \emph{instantiate} the autonomous, semi-independent module dynamics that characterize biological oscillator assemblies. This is a hypothesis about dynamical structure, not a theorem; intermediate architectures (asynchronous digital, neuromorphic with local clocks) occupy a spectrum between the extremes we describe.

\paragraph{Unclocked architectures.} Biological neural networks, and potentially neuromorphic hardware designed on similar principles, operate without global synchronization. Each module evolves on its own trajectory; the phase space is the \emph{actual} dynamical space of the system, not a represented abstraction. The dimensionality is ontologically realized in the substrate.

Such systems pay the coherence time penalty derived in this paper: coordination across $M$ modules requires waiting for rare alignment events, with time scaling exponentially in $M$. But they gain access to genuinely high-dimensional dynamics---the very property that prior work identified as constitutive of intelligence~\cite{todd2025intelligence}.

\paragraph{A design hypothesis.} This suggests a potential trade-off for artificial systems:
\begin{itemize}
\item \textbf{Clocked}: Fast, deterministic, scalable in parallelism---but potentially constrained in dynamical autonomy. The substrate may not support the spontaneous coordination dynamics modeled here.
\item \textbf{Unclocked}: Capable of genuinely autonomous, high-dimensional dynamics---but subject to coordination bottlenecks that scale exponentially with integration depth.
\end{itemize}

If high-dimensional coherent dynamics prove to be not merely \emph{correlated with} but \emph{constitutive of} certain cognitive capacities, then clocked architectures may face constraints that unclocked ones do not. This remains speculative; the empirical question is whether the coordination-time framework applies to engineered systems and, if so, whether it identifies real performance boundaries.

This framing may inform the design of neuromorphic systems, spiking neural networks, and other non-von Neumann architectures. The coherence time formula (Eq.~\ref{eq:coh-time}) provides quantitative guidance: to achieve update rates comparable to biological cognition, such systems must either accept modest coordination depth $M$, engineer high baseline coherence $r$, or develop hierarchical architectures that localize coordination requirements.

\paragraph{The cerebellum as biological precedent.} The mammalian cerebellum exemplifies this hierarchical solution. Operating with low coordination depth ($M \sim 3$--5 modules), tight coupling ($r \sim 0.8$), and stereotyped primitives, the cerebellum achieves commit times of 10--20 ms---fast enough to override slow, high-dimensional cortical dynamics when rapid action is required. Skilled movement, postural correction, and overlearned sequences route through this low-$M$ pathway precisely because cortical coordination cannot keep pace.

Notably, cerebellar processing is largely \emph{unconscious}. We do not experience the moment-to-moment adjustments during a golf swing or the postural corrections that keep us upright. If phenomenal experience is associated with high-dimensional pre-commit dynamics (as suggested in \S4.4), then low-dimensional cerebellar override would be expected to fall outside awareness. The architecture that solves the speed problem---collapsing to low $M$ for fast commits---simultaneously removes processing from the domain of conscious experience. This is not a bug but a feature: rapid action requires precisely the dimensional compression that excludes phenomenology.

\section{Conclusion}

We have established that coherence time---the time required for distributed oscillatory assemblies to achieve sufficient phase alignment for irreversible state registration---bounds the rate of neural processes associated with cognition. The lower bound
\[
\tau_{\mathrm{eff}} \gtrsim \max\left[\tau_{\mathrm{QSL}},\, \tau_{\mathrm{SNR}},\, \tau_{\mathrm{coh}},\, \tau_{\mathrm{power}}\right]
\]
with coherence time
\[
\tau_{\mathrm{coh}} = \frac{1}{\lambda_{\mathrm{attempt}}}\,p_1(\varepsilon, \kappa(r))^{-\alpha(M-1)}
\]
unifies quantum, noise, and coordination limits to explain temporal resolution in biological processing systems.

We have situated this framework explicitly within the hierarchy of physical limits on computation: the Landauer bound constrains energy per irreversible operation; the Margolus-Levitin and Mandelstam-Tamm bounds constrain the speed of quantum state evolution; the energy-time uncertainty relation underlies both. In biological neural systems, these quantum and thermodynamic bounds are typically non-binding. The dominant constraint is the time required for multi-module coordination---a classical phenomenon arising from rare-event statistics on product manifolds.

Key findings:
\begin{enumerate}
\item \textbf{Speed-flexibility trade-off}: Increasing $M$ exponentially slows commits but expands combinatorial flexibility. Increasing $r$ speeds commits but restricts dynamics to low-dimensional manifolds.

\item \textbf{Quantum limits non-binding}: Neural coherence times exceed quantum speed limits by eleven orders of magnitude. Coordination, not quantum switching, limits processing rate.

\item \textbf{Generality}: While neural systems provide the empirical testbed, the framework applies to any biological oscillator assembly with coordinated phase relations and sparse commits.

\item \textbf{Quantitative predictions}: Visual binding windows (30--50 ms), metabolic scaling of temporal acuity, alpha entrainment linearity, dual-task dissociations.
\end{enumerate}

The framework also has potential clinical implications. Pathologies that degrade inter-module coupling should reduce effective coordination capacity ($\alpha$ or $M$) before individual components fail overtly. In neurodegenerative disease, the coherence time framework predicts that coordination-level signatures---changes in the scaling of inter-commit intervals with task complexity---should precede overt motor or cognitive deficits, because degradation of inter-module coupling affects the exponent $\alpha(M-1)$ before individual neurons are lost. Motor neuron disease offers one example: cortical hyperexcitability in ALS is mediated by dysfunction of distinct interneuronal circuits rather than uniform degeneration~\cite{pavey2024}, suggesting that coordination-level biomarkers derived from the coherence time framework may complement single-neuron measures.

The central thesis: biological systems trade speed for flexibility by moving on the coherence-dimension frontier, operating in a regime where quantum and thermodynamic limits permit far faster processing than coordination constraints allow.

\section*{Data accessibility}

This is a theoretical framework paper. Kuramoto simulation code and parameter estimation workflows are available at \url{https://github.com/todd866/coherence-time}.

\section*{Competing interests}

The author declares no competing interests.

\section*{Funding}

No external funding supported this work.

\section*{Acknowledgments}

The author thanks Abir Igamberdiev and anonymous reviewers at BioSystems for constructive feedback.

\begin{thebibliography}{99}

\bibitem{miller2018}
Miller, E.K., Lundqvist, M., Bastos, A.M. (2018). Working memory 2.0. \textit{Neuron}, 100(2), 463--475.

\bibitem{lundqvist2016}
Lundqvist, M., Rose, J., Herman, P., Brincat, S.L., Buschman, T.J., Miller, E.K. (2016). Gamma and beta bursts underlie working memory. \textit{Neuron}, 90(1), 152--164.

\bibitem{pinotsis2023}
Pinotsis, D.A., Fridman, G., Miller, E.K. (2023). Cytoelectric coupling: Electric fields sculpt neural activity and ``tune'' the brain's infrastructure. \textit{Progress in Neurobiology}, 226, 102465.

\bibitem{fries2015}
Fries, P. (2015). Rhythms for cognition: communication through coherence. \textit{Neuron}, 88(1), 220--235.

\bibitem{singer1999}
Singer, W. (1999). Neuronal synchrony: a versatile code for the definition of relations? \textit{Neuron}, 24(1), 49--65.

\bibitem{lisman2013}
Lisman, J.E., Jensen, O. (2013). The theta-gamma neural code. \textit{Neuron}, 77(6), 1002--1016.

\bibitem{womelsdorf2007}
Womelsdorf, T., et al. (2007). Modulation of neuronal interactions through neuronal synchronization. \textit{Science}, 316(5831), 1609--1612.

\bibitem{engel2001}
Engel, A.K., Fries, P., Singer, W. (2001). Dynamic predictions: oscillations and synchrony in top-down processing. \textit{Nat. Rev. Neurosci.}, 2(10), 704--716.

\bibitem{varela2001}
Varela, F., et al. (2001). The brainweb: phase synchronization and large-scale integration. \textit{Nat. Rev. Neurosci.}, 2(4), 229--239.

\bibitem{fries2005}
Fries, P. (2005). A mechanism for cognitive dynamics: neuronal communication through neuronal coherence. \textit{Trends Cogn. Sci.}, 9(10), 474--480.

\bibitem{todd2025intelligence}
Todd, I. (2026). Intelligence as high-dimensional coherence: The observable dimensionality bound and computational tractability. \textit{BioSystems}, 260, 105704.

\bibitem{landauer1961}
Landauer, R. (1961). Irreversibility and heat generation in the computing process. \textit{IBM J. Res. Dev.}, 5(3), 183--191.

\bibitem{bormashenko2024}
Bormashenko, E. (2024). Landauer bound in the context of minimal physical principles: Meaning, experimental verification, controversies and perspectives. \textit{Entropy}, 26(5), 423.

\bibitem{mandelstam1945}
Mandelstam, L., Tamm, I. (1945). The uncertainty relation between energy and time in non-relativistic quantum mechanics. \textit{J. Phys. (USSR)}, 9, 249--254.

\bibitem{margolus1998}
Margolus, N., Levitin, L.B. (1998). The maximum speed of dynamical evolution. \textit{Physica D}, 120(1--2), 188--195.

\bibitem{khrennikov2021}
Khrennikov, A. (2021). Quantum-like model for unconscious--conscious interaction and emotional coloring of perceptions and other conscious experiences. \textit{BioSystems}, 208, 104471.

\bibitem{igamberdiev2025constraints}
Igamberdiev, A.U., Shklovskiy-Kordi, N.E. (2025). Physical limits of natural computation as the biological constraints of morphogenesis, evolution, and consciousness: On the 100th anniversary of Efim Liberman (1925--2011). \textit{BioSystems}, 251, 105451.

\bibitem{igamberdiev2026autopoietic}
Igamberdiev, A.U. (2026). Internal quantum constraints of natural computation in autopoietic systems. \textit{BioSystems}, 261, 105707.

\bibitem{khrennikov2025bridge}
Khrennikov, A., Iriki, A., Basieva, I. (2025). Constructing a bridge between functioning of oscillatory neuronal networks and quantum-like cognition along with quantum-inspired computation and AI. \textit{BioSystems}, 257, 105573.

\bibitem{poor1994}
Poor, H.V. (1994). \textit{An Introduction to Signal Detection and Estimation} (2nd ed.). Springer-Verlag.

\bibitem{kac1947}
Kac, M. (1947). On the notion of recurrence in discrete stochastic processes. \textit{Bull. Amer. Math. Soc.}, 53(10), 1002--1010.

\bibitem{mardia2000}
Mardia, K.V., Jupp, P.E. (2000). \textit{Directional Statistics}. John Wiley \& Sons.

\bibitem{stetson2007}
Stetson, C., Fiesta, M.P., Eagleman, D.M. (2007). Does time really slow down during a frightening event? \textit{PLoS ONE}, 2(12), e1295.

\bibitem{samaha2015}
Samaha, J., Postle, B.R. (2015). The speed of alpha-band oscillations predicts the temporal resolution of visual perception. \textit{Current Biology}, 25(22), 2985--2990.

\bibitem{cecere2015}
Cecere, R., Rees, G., Romei, V. (2015). Individual differences in alpha frequency drive crossmodal illusory perception. \textit{Current Biology}, 25(2), 231--235.

\bibitem{healy2013}
Healy, K., McNally, L., Ruxton, G.D., Cooper, N., Jackson, A.L. (2013). Metabolic rate and body size are linked with perception of temporal information. \textit{Anim. Behav.}, 86(4), 685--696.

\bibitem{takens1981}
Takens, F. (1981). Detecting strange attractors in turbulence. In: Rand, D., Young, L.S. (eds) \textit{Dynamical Systems and Turbulence}, Lecture Notes in Mathematics, vol 898. Springer, pp. 366--381.

\bibitem{braitenberg1958}
Braitenberg, V., Atwood, R.P. (1958). Morphological observations on the cerebellar cortex. \textit{Journal of Comparative Neurology}, 109, 1--33.

\bibitem{llinas1982}
Llin\'{a}s, R. (1982). General discussion: Radial connectivity in the cerebellar cortex. In: \textit{The Cerebellum---New Vistas} (pp. 189--194). Springer.

\bibitem{wyatt2005}
Wyatt, K.D., et al. (2005). Speed limits in the cerebellum: constraints from myelinated and unmyelinated parallel fibers. \textit{European Journal of Neuroscience}, 21(8), 2285--2290.

\bibitem{tononi2016}
Tononi, G., Boly, M., Massimini, M., Koch, C. (2016). Integrated information theory: from consciousness to its physical substrate. \textit{Nat. Rev. Neurosci.}, 17(7), 450--461.

\bibitem{friston2010}
Friston, K. (2010). The free-energy principle: a unified brain theory? \textit{Nat. Rev. Neurosci.}, 11(2), 127--138.

\bibitem{anastassiou2011}
Anastassiou, C.A., Perin, R., Markram, H., Koch, C. (2011). Ephaptic coupling of cortical neurons. \textit{Nat. Neurosci.}, 14(2), 217--223.

\bibitem{beggs2003}
Beggs, J.M., Plenz, D. (2003). Neuronal avalanches in neocortical circuits. \textit{J. Neurosci.}, 23(35), 11167--11177.

\bibitem{shew2013}
Shew, W.L., Plenz, D. (2013). The functional benefits of criticality in the cortex. \textit{Neuroscientist}, 19(1), 88--100.

\bibitem{han2018}
Han, K.S., Guo, C., Chen, C.H., Witter, L., Osorno, T., Regehr, W.G. (2018). Ephaptic coupling promotes synchronous firing of cerebellar Purkinje cells. \textit{Neuron}, 100(3), 564--578.

\bibitem{pavey2024}
Pavey, N., Hannaford, A., van den Bos, M., Kiernan, M.C., Menon, P., Vucic, S. (2024). Distinct neuronal circuits mediate cortical hyperexcitability in amyotrophic lateral sclerosis. \textit{Brain}, 147(7), 2344--2356.

\bibitem{kuramoto1984}
Kuramoto, Y. (1984). \textit{Chemical Oscillations, Waves, and Turbulence}. Springer-Verlag, Berlin.

\bibitem{strogatz2000}
Strogatz, S.H. (2000). From Kuramoto to Crawford: exploring the onset of synchronization in populations of coupled oscillators. \textit{Physica D}, 143(1-4), 1--20.

\bibitem{acebron2005}
Acebr\'{o}n, J.A., Bonilla, L.L., P\'{e}rez Vicente, C.J., Ritort, F., and Spigler, R. (2005). The Kuramoto model: A simple paradigm for synchronization phenomena. \textit{Reviews of Modern Physics}, 77(1), 137--185.

\bibitem{ott2008}
Ott, E. and Antonsen, T.M. (2008). Low dimensional behavior of large systems of globally coupled oscillators. \textit{Chaos}, 18(3), 037113.

\bibitem{cunningham2014}
Cunningham, J.P. and Yu, B.M. (2014). Dimensionality reduction for large-scale neural recordings. \textit{Nature Neuroscience}, 17(11), 1500--1509.

\bibitem{stringer2019}
Stringer, C., Pachitariu, M., Steinmetz, N., Reddy, C.B., Carandini, M., and Harris, K.D. (2019). Spontaneous behaviors drive multidimensional, brainwide activity. \textit{Science}, 364(6437), eaav7893.

\bibitem{todd2025biosystems}
Todd, I. (2025). The limits of falsifiability in continuously adapting biological systems. \textit{BioSystems}, 258, 105608.

\bibitem{todd2025demon}
Todd, I. (2025). Timing inaccessibility in biological Maxwell demons. \textit{BioSystems}, 258, 105632.

\bibitem{raichle2001}
Raichle, M.E., MacLeod, A.M., Snyder, A.Z., Powers, W.J., Gusnard, D.A., and Shulman, G.L. (2001). A default mode of brain function. \textit{Proceedings of the National Academy of Sciences}, 98(2), 676--682.

\bibitem{wolpert1998}
Wolpert, D.M., Miall, R.C., and Kawato, M. (1998). Internal models in the cerebellum. \textit{Trends in Cognitive Sciences}, 2(9), 338--347.

\bibitem{vanrullen2016}
VanRullen, R. (2016). Perceptual cycles. \textit{Trends in Cognitive Sciences}, 20(10), 723--735.

\bibitem{helfrich2014}
Helfrich, R.F., Schneider, T.R., Rach, S., Trautmann-Lengsfeld, S.A., Engel, A.K., and Herrmann, C.S. (2014). Entrainment of brain oscillations by transcranial alternating current stimulation. \textit{Current Biology}, 24(3), 333--339.

\bibitem{amari2016}
Amari, S. (2016). \textit{Information Geometry and Its Applications}. Springer, Tokyo.

\bibitem{ashby1956}
Ashby, W.R. (1956). \textit{An Introduction to Cybernetics}. Chapman \& Hall, London.

\bibitem{gabriel1996}
Gabriel, S., Lau, R.W., Gabriel, C. (1996). The dielectric properties of biological tissues: III. Parametric models for the dielectric spectrum of tissues. \textit{Phys. Med. Biol.}, 41(11), 2271--2293.

\bibitem{subramanian2022}
Subramanian, M., Chiang, C.-C., Couturier, N.H., Durand, D.M. (2022). Theta waves, neural spikes and seizures can propagate by ephaptic coupling in vivo. \textit{Exp. Neurol.}, 354, 114109.

\end{thebibliography}

\end{document}
